\section{Methods}

\subsection{Dataset and activity-cliff definition}
\TODO{ChEMBL target CHEMBL220 (human AChE); pChEMBL aggregation; similarity and
  potency thresholds; 1{,}911 similar pairs; cliff/control definition.}
\NOTE{Source: results/01\_*, 02\_pairs\_CHEMBL220.csv, 02\_cliff\_similarity\_audit.csv}

\subsection{Stereochemical cliffs are an artefact of stereo-blind fingerprints}
\TODO{Enantiomer pairs reach Tanimoto 1.0 under ECFP and are counted as cliffs
  that no 2D method can resolve by construction, while 3D methods resolve them
  trivially. Report how many, and how they are handled here.}
\NOTE{This biases cliff benchmarks toward 3D methods -- the opposite direction to
  the subset artefact in Section~\ref{sec:subset}. Having both measured in one
  study is the point.}

\subsection{Splits}
\TODO{Connected-component splits over the similarity graph, seed 20260902, five
  CV folds, sealed test set of 54 ligands. State that the test split has not been
  scored by any model and when it will be.}

\subsection{The three framings}
\TODO{A: per-ligand $R^2$ on pChEMBL. B (primary): $R^2$ on the signed pair
  difference, antisymmetric rows. C: cliff/control classification, ROC AUC,
  order-invariant pair features $[d\_a+d\_b,\,|d\_a-d\_b|]$.}

\subsection{Descriptor rungs}
\TODO{Rung 1: binary ECFP4, and count ECFP4 + RDKit2D -- both, because the winner
  depends on the framing. Rung 2: pose geometry on 4EY7 plus bridging features.
  Rung 2e: median over eight receptor conformations. Rung 3: single-snapshot
  MM/GBSA, energy and per-residue terms only. Rung 4: MD descriptors, below.}

\subsection{Molecular dynamics}
\TODO{OpenFF Sage 2.3.0 + amber14SB + TIP3P; system build; equilibration; 25 ns
  production (45 ns for the first 55 runs); GPU-resident mdrun; 217 ligands.}
\NOTE{Receptor preparation matters and belongs here: naive loop bridging in 4EY7
  created fabricated bonds worth 90\% of the bond energy; MD uses the capped
  4ey7\_rec\_md.pdb. Source: results/12b\_receptor\_chain\_breaks.md}

\subsection{MD-derived descriptors}
\TODO{33 descriptors, three design classes -- persistence, fluctuation, event
  count -- deliberately not frame-averaged rung-2 quantities. Table of
  definitions. 14 gorge-lining residues; CAS = Trp86, PAS = Trp286, gate = Tyr337.}
\NOTE{Of the three classes only persistence survives the replicate filter: 8 of
  the 11 survivors are contact occupancies. Report that here or in Results, but
  report it -- it is the actionable guidance for anyone building MD descriptors.}

\subsection{Periodic boundary handling}
\TODO{Bond-based unwrap, nojump, minimum image, Kabsch superposition. State that
  an earlier version was wrong and that 19 of 30 descriptors changed when it was
  fixed, with the A/B test as evidence the QC catches real defects.}
\NOTE{Source: results/16b\_rung4\_pbc\_fix.md}

\subsection{Replicate pilot and the descriptor filter}
\TODO{Reliability $\sigma^2\_{\mathrm{lig}}/(\sigma^2\_{\mathrm{lig}}+2\sigma^2\_{\mathrm{rep}})$;
  12 replicate runs over 5 ligands; 81 of 160 descriptor$\times$window combinations
  are seed-dominated; the filter keeps 11 at the 20 ns window.}
\NOTE{Source: results/18\_replicate\_pilot.md and 18\_replicate\_pilot.csv}

\subsection{Model suite and evaluation}
\TODO{Six regression and five classification families; preprocessing fitted inside
  each fold; best-of-family reported; bootstrap over units with the algorithm --
  and, for rung 1, the feature set -- re-selected inside every replicate.}
\NOTE{State the bootstrap's known bias: resampling units and then running the
  frozen folds puts duplicates on both sides of the split and inflates the levels.
  It applies identically to both sides, so the gap is the quantity to read.
  Bootstrap medians must not be quoted as performance.}

\subsection{Statistical power}
\TODO{Synthetic descriptor matrices through the same pipeline; power against the
  rung-1 baseline with the baseline re-simulated each replicate; measurement-noise
  ceiling from repeat assays.}
